\documentclass[header={Lecture 3 Exam Study Notes}]{study-note}

\begin{document}

\StudyTitleBlock{03}{Lecture 3: Exam Study Notes}{Seminorms, Fr\'echet spaces, Hahn--Banach theorem, and separation}

\vspace{0.8em}

\begin{focusbox}
This lecture is exam-critical. It contains the first statement of the Hahn--Banach theorem, the distance generated by separating seminorms, and the first separation consequences used later for weak convergence and duality.
\end{focusbox}

\tableofcontents

\section{What To Learn}

\begin{focusbox}
You should know:
\begin{itemize}
\item what a seminorm is and why one seminorm alone need not define a metric;
\item how a \emph{separating sequence} of seminorms generates a metric;
\item the real and complex Hahn--Banach extension theorems;
\item topological Hahn--Banach and its main consequences;
\item separation of a point from a closed subspace;
\item dual characterization of the norm;
\item why separability of $X^*$ implies separability of $X$.
\end{itemize}
\end{focusbox}

\section{Seminorms And Fr\'echet Spaces}

\begin{defbox}{Seminorm}
Let $X$ be a vector space over $\K$, where $\K=\R$ or $\C$. A map $p:X\to\R$ is a \textbf{seminorm} if:
\begin{align*}
\text{(SN1)}\quad & p(x)\ge 0,
\\
\text{(SN2)}\quad & p(\lambda x)=\abs{\lambda}p(x) \qquad \forall \lambda\in\K,
\\
\text{(SN3)}\quad & p(x+y)\le p(x)+p(y).
\end{align*}
Unlike a norm, one may have $p(x)=0$ for some $x\neq 0$.
\end{defbox}

\begin{defbox}{Separating sequence of seminorms}
A sequence $(p_k)_{k\in\N}$ of seminorms on $X$ is \textbf{separating} if for every nonzero $x\in X$ there exists $k$ such that $p_k(x)>0$.
\end{defbox}

\begin{thmbox}{Theorem 5.1: A separating sequence of seminorms generates a metric}
\thmFrechetMetric
\end{thmbox}

\begin{warningbox}{Common trap}
Separation of the seminorms is not needed for the triangle inequality of the generated distance. It is needed only for $d(x,y)=0\Rightarrow x=y$.
\end{warningbox}

\begin{prooflevelbox}{SKETCH}
Know why each axiom holds. The triangle inequality uses subadditivity of $t/(1+t)$. Separation is needed only for $d(x,y)=0\Rightarrow x=y$.
\end{prooflevelbox}

\paragraph{Proof idea.}
Non-negativity and symmetry are clear. If $d(x,y)=0$, then $p_k(x-y)=0$ for all $k$, and separation gives $x=y$. The triangle inequality follows because the function
\[
t\mapsto \frac{t}{1+t}
\]
is increasing and subadditive on $[0,\infty)$.

\section{Hahn--Banach Extension}

\begin{defbox}{Sublinear functional}
Let $X$ be a \emph{real} vector space. A map $p:X\to\R$ is \textbf{sublinear} if
\[
p(x+y)\le p(x)+p(y),
\qquad
p(t x)=t p(x) \text{ for every } t\ge 0.
\]
\end{defbox}

\begin{thmbox}{Theorem 6.1: Real algebraic Hahn--Banach theorem}
\thmHahnBanachReal
\end{thmbox}

\begin{prooflevelbox}{SKETCH}
Know the one-step extension (choose $F(x_0)$ inside the interval forced by sublinearity) and that Zorn's lemma finishes the argument.
\end{prooflevelbox}

\paragraph{Proof idea.}
Extend first from $V$ to a one-dimensional enlargement $V\oplus \operatorname{span}\{x_0\}$ by choosing the value of $F(x_0)$ inside an interval forced by sublinearity. Then apply Zorn's lemma to the partially ordered family of all admissible extensions.

\begin{thmbox}{Theorem 6.1': Complex Hahn--Banach theorem}
\thmHahnBanachComplex
\end{thmbox}

\begin{prooflevelbox}{SKETCH}
Apply real Hahn--Banach to $\operatorname{Re}f$, then reconstruct the complex extension via $F(x)=U(x)-iU(ix)$.
\end{prooflevelbox}

\paragraph{Proof idea.}
Apply the real Hahn--Banach theorem to the real part of $f$ on the underlying real vector space, then reconstruct the complex-linear extension by
\[
F(x)=U(x)-iU(ix).
\]

\begin{thmbox}{Theorem 6.2: Topological Hahn--Banach theorem}
\thmHahnBanachNormed
\end{thmbox}

\begin{warningbox}{Common trap}
Hahn--Banach extends continuous functionals with the same norm, but it does not say the extension is unique.
\end{warningbox}

\begin{prooflevelbox}{FULL}
Know how to choose $p(x)=\norm{f}\norm{x}$ and apply the algebraic theorem. Verify that the extension satisfies $\abs{F(x)}\le\norm{f}\norm{x}$ and hence is bounded with the same norm.
\end{prooflevelbox}

\paragraph{Proof.}
Apply the algebraic Hahn--Banach theorem with the sublinear functional
\[
p(x):=\norm{f}\norm{x}.
\]
The extension automatically satisfies $\abs{F(x)}\le \norm{f}\norm{x}$, hence is bounded. Equality of norms follows because $F$ extends $f$.

\section{Separation And Consequences}

\begin{thmbox}{Theorem 6.3: Separation of a point from a closed subspace}
\thmHahnBanachSeparating
\end{thmbox}

\begin{prooflevelbox}{SKETCH}
Define a functional on $Y+\operatorname{span}\{x_0\}$ by $y+tx_0\mapsto t\,\operatorname{dist}(x_0,Y)$, verify boundedness, then extend via Theorem 6.2.
\end{prooflevelbox}

\paragraph{Proof idea.}
Define a linear functional first on $Y+\operatorname{span}\{x_0\}$ by sending $y+tx_0$ to $t\,\operatorname{dist}(x_0,Y)$ and verify it is bounded by the norm. Then extend it with Theorem~\ref{thm:hahn-banach-normed}.

\begin{thmbox}{Theorem 6.4: Separation of a nonzero element from the origin}
\thmDualSeparates
\end{thmbox}

\begin{prooflevelbox}{STATEMENT}
This is the special case $Y=\{0\}$ of Theorem 6.3.
\end{prooflevelbox}

\paragraph{Reason.}
Apply Theorem~\ref{thm:hahn-banach-separating} with $Y=\{0\}$. If $x_n\rightharpoonup x$ and $x_n\rightharpoonup y$, then $\phi(x-y)=0$ for all $\phi\in X^*$, hence $x=y$.

\begin{thmbox}{Theorem 6.5: Dual characterization of the norm}
\thmNormViaDual
\end{thmbox}

\begin{prooflevelbox}{FULL}
One inequality: $\abs{\phi(x)}\le\norm{\phi}\norm{x}$. Reverse: apply Theorem 6.3 to $\{0\}$ to find a norm-one functional attaining $\norm{x}$.
\end{prooflevelbox}

\paragraph{Proof.}
One inequality is immediate from
\[
\abs{\phi(x)}\le \norm{\phi}\norm{x}.
\]
For the reverse inequality, apply Theorem~\ref{thm:hahn-banach-separating} to the one-point subspace $\{0\}$ to get a norm-one functional attaining $\norm{x}$ on $x$.

\begin{thmbox}{Theorem 6.6: If $X^*$ is separable, then $X$ is separable}
\thmDualSeparableImpliesSeparable
\end{thmbox}

\begin{prooflevelbox}{SKETCH}
Choose dense $(\phi_n)$ in $X^*$ and norm-approaching $(x_n)$ in $X$. Show the closed span of $(x_n)$ is all of $X$; otherwise Hahn--Banach contradicts density.
\end{prooflevelbox}

\paragraph{Proof idea.}
Choose a dense sequence $(\phi_n)$ in $X^*$ and points $(x_n)$ that almost realize their norms. The closed span of $(x_n)$ is then all of $X$; otherwise Hahn--Banach separation would contradict density of $(\phi_n)$.

\section{Exam Orientation}

\begin{warningbox}{What exam questions this lecture supports}
This lecture directly feeds into oral exam topics:
\begin{itemize}
\item Hahn--Banach theorem: formulation and proof idea;
\item consequences of Hahn--Banach: separation, weak convergence, dual norm formula, canonical embedding;
\item definition of dual space and why there are enough functionals.
\end{itemize}
\end{warningbox}

\section{What To Memorize Precisely}

\begin{focusbox}
Be able to state without hesitation:
\begin{enumerate}
\item the definitions of seminorm, separating sequence, sublinear functional, Fr\'echet space, dual space;
\item Theorems \textbf{5.1}, \textbf{6.1}, \textbf{6.1'}, \textbf{6.2}, \textbf{6.3}, \textbf{6.4}, \textbf{6.5}, \textbf{6.6};
\item why Hahn--Banach is the reason weak convergence even makes sense;
\item the norm formula
\[
\norm{x}=\sup_{\norm{\phi}\le 1}\abs{\phi(x)};
\]
\item the role of Zorn's lemma in the algebraic proof.
\end{enumerate}
\end{focusbox}

\end{document}
